\chapter{Introduction}
\label{ch:introduction}

\epigraph{When water flashes to steam, it will always attempt to increase to
roughly $1{,}700$ times its original volume.}{\textit{The National Board of
Boiler and Pressure Vessel Inspectors} \citep{nationalboard-steam}}

\section{Background and Motivation}
\label{sec:motivation}

For most of computing history, the only ``language'' a machine could be made to
understand with any depth was source code, and even that understanding was
trapped inside individual tools. An editor that knew how to complete a Python
symbol knew nothing about Rust; a tool that could navigate Java could not hover
over TypeScript. Each pair of an editor and a programming language demanded its
own integration, so that supporting $M$ editors across $N$ languages implied, in
the worst case, $M \times N$ separate implementations of the same handful of
ideas: completion, go-to-definition, diagnostics, hover.

In 2016, Microsoft, Red Hat, and Codenvy published the \emph{Language Server
Protocol} (LSP), a JSON-RPC contract that let a single \emph{language server}
serve any number of editors \citep{lsp-redhat-2016,lsp-official}. The protocol's
own documentation states the design intent plainly: ``a single Language Server
can be re-used in multiple development tools, which in turn can support multiple
languages with minimal effort'' \citep{lsp-official}. The combinatorial effect
is the move from $M \times N$ to $M + N$: $M$ editor adapters plus $N$ servers
\citep{lsp-mxn-fcc}. LSP did not make editors smarter; it \emph{decoupled} the
locus of language intelligence from the locus of display, and in doing so turned
``understanding a language'' into a reusable, networked service.

A decade later, the same decoupling has been performed twice more, one layer up
the stack, in response to the arrival of capable language models. In November
2024 Anthropic open-sourced the \emph{Model Context Protocol} (MCP), ``a new
standard for connecting AI assistants to the systems where data lives''
\citep{mcp-anthropic-2024}; its documentation likens it to ``a USB-C port for AI
applications'' \citep{mcp-intro}, and---tellingly for this thesis---its
specification records that ``MCP takes some inspiration from the Language Server
Protocol'' \citep{mcp-spec-2025}. In April 2025 Google published the
\emph{Agent-to-Agent} protocol (A2A), which standardises how autonomous agents
``communicate with each other, securely exchange information, and coordinate
actions'' and which Google explicitly frames as complementary to MCP: ``A2A is an
open protocol that complements Anthropic's Model Context Protocol (MCP), which
provides helpful tools and context to agents'' \citep{a2a-google-2025}. Where LSP
decoupled language intelligence from the editor, MCP decouples context and tools
from the model, and A2A decouples agent capability from the framework that hosts
it.

These three protocols are no longer merely conceptually adjacent; they are
converging institutionally. A2A was donated to the Linux Foundation in June 2025
\citep{a2a-lf-donation,lf-a2a-press}, and in December 2025 MCP was contributed to
the newly formed \emph{Agentic AI Foundation} under the same neutral umbrella,
alongside further interoperability projects \citep{mcp-aaif-2025,lf-aaif-press}.
The substrate on which the next decade of machine-mediated language will run is,
visibly, being poured.

Yet a gap remains, and it is the gap this thesis addresses. The protocols above
expand \emph{reach}---which tools a model can call, which agents an agent can
summon---without expanding \emph{accountability}. They standardise how messages
flow, not whether the behaviour those messages induce conforms to any model of
how it ought to unfold. This is precisely the question that an entire scientific
discipline already answers for business processes.

That discipline is \emph{process mining}, founded by Wil van der Aalst, whose
central idea is ``to discover, monitor and improve real processes \dots\ by
extracting knowledge from event logs readily available in today's
\dots\ systems'' \citep{pm-manifesto-2011,vdaalst-book-2016}. Process mining
already possesses a mature theory of \emph{conformance checking}: given a
recorded event log and a model of intended behaviour, it measures where reality
and model agree and where they diverge \citep{conformance-review-2020}. Van der
Aalst has recently argued that object-centric process mining is ``the missing
link connecting data and processes'' and the prerequisite for trustworthy
generative, predictive, and prescriptive AI---``no AI without PI'' (process
intelligence) \citep{noai-nopi-2025}. This thesis takes that argument to its
conclusion across \emph{all} language, not only business process.

\section{Problem Statement}
\label{sec:problem}

Three problems motivate the work.

\paragraph{P1 --- The narrowness of machine language.} The infrastructure for
deep, tool-mediated language understanding (LSP and its descendants) was built
for source code. Natural language, legal and regulatory text, mathematics,
biological sequence, music, and business process each have rich structure and
rich notions of correctness, yet none enjoys the editor-grade, server-mediated
intelligence that code does. The generalisation of LSP beyond code is asserted in
isolated projects---grammar servers such as LTeX \citep{ltex}, data-format
servers such as the YAML language server \citep{yamlls}---but lacks a unifying
account of \emph{when} a domain is ``a language'' that a language server can
serve.

\paragraph{P2 --- Reach without accountability.} MCP and A2A multiply what
agents can do and whom they can ask, but they carry no native notion of
conformance. An agent may call a tool, receive a result, and act, with no
standing surface that asks whether the trajectory conformed to a declared model
of lawful behaviour. As agentic systems scale---Gartner forecasts that a large
fraction of agentic-AI projects will be cancelled by 2027 for, among other
reasons, ``inadequate risk controls'' \citep{gartner-agentic-2025}---the absence
of a conformance surface becomes a structural, not cosmetic, deficiency.

\paragraph{P3 --- Fragmentation versus fusion.} The interoperability landscape is
often narrated as a ``protocol war'' among MCP, A2A, and alternatives such as
ACP, ANP, and AGNTCY \citep{acp-a2a-merge,anp-github,agntcy-cisco}. This thesis
contends the opposite: that LSP, MCP, and A2A are not competitors but
\emph{strata} of one substrate, and that their fusion---rather than the dominance
of any one---is the event of consequence.

\section{The Central Thesis}
\label{sec:thesis-statement}

\begin{hypothesis}[The phase-transition thesis]
\label{hyp:central}
When the Language Server Protocol is generalised from source code to all forms of
language and fused with the Model Context Protocol and the Agent-to-Agent
protocol, the addressable surface of machine-mediated language undergoes a
\emph{phase transition}: a discontinuous, latent-heat-driven expansion---on the
order of the ${\sim}1{,}600$--$1{,}700\times$ expansion of liquid water into
steam at its boiling point---in which language ceases to be a thing models
\emph{read} and becomes a medium agents \emph{inhabit, observe, and are held to
account within}. The enabling abstraction for trustworthy expansion is
conformance checking: every language, once recorded as an event log, acquires a
conformance surface.
\end{hypothesis}

The phase-transition framing is developed formally in \cref{ch:framework}. Its
force is twofold. First, it is \emph{quantitative}: phase transitions are not
gradual: water at $99\,^{\circ}\mathrm{C}$ and water at $100\,^{\circ}\mathrm{C}$
differ by a degree, but liquid water and steam differ in volume by roughly three
orders of magnitude \citep{phys-illinois-steam,nationalboard-steam}. Second, it
is \emph{energetic}: the expansion is paid for, all at once, by a large input
absorbed at constant temperature---the latent heat of vaporisation, about
$2{,}256\ \mathrm{kJ/kg}$ for water \citep{latentheat-openstax}. The analogue of
that latent heat, we argue, is \emph{standardisation}: the unglamorous,
discontinuous investment of agreeing on a protocol, which yields no visible
progress while it is being paid and an order-of-magnitude expansion once it is.

\section{Research Questions}
\label{sec:rqs}

\rq{1}{What structural property do LSP, MCP, and A2A share, and in what sense is
their composition a single substrate rather than three protocols?}

\rq{2}{Under what conditions does generalising LSP beyond source code to an
arbitrary domain become coherent---that is, what must a domain provide in order
to be served by a ``language server''?}

\rq{3}{How can van der Aalst's process-mining paradigm---event logs, conformance
checking, and object-centric event logs (OCEL)---be applied to all such language,
and what does it add that the protocols alone do not provide?}

\rq{4}{What artifact design realises a law-state runtime that brings conformance
to all language while remaining trustworthy---read-only with respect to the
artifacts it judges, receipt-bound in its claims, and three-valued in its
verdicts?}

\rq{5}{Is the convergence of these protocols best understood as incremental
integration or as a phase transition, and what does the metaphor, taken
seriously, predict for the year 2030?}

\section{Objectives, Scope, and Contributions}
\label{sec:contributions}

The objective is a coherent, evidence-grounded, and falsifiable account of
protocol convergence as a phase transition, instantiated in a concrete artifact.
The scope is conceptual and design-scientific: we build and analyse an artifact
(\lspmax) and reason about a 2030 horizon; we do \emph{not} claim a completed
empirical deployment across all six language domains, and we are explicit
(\cref{ch:discussion}) about which claims are \admitted, which are \candidate,
and which remain \openstat. The contributions are:

\begin{itemize}[leftmargin=2.2em]
  \item[\textbf{C1}] A unifying account (\cref{ch:background,ch:framework}) of
  LSP, MCP, and A2A as three instances of one decoupling, each turning an
  $M \times N$ integration burden into $M + N$ at a different layer, and
  composing into a single substrate.
  \item[\textbf{C2}] A criterion (\cref{ch:all-language}) for when a domain is
  ``a language'' amenable to a language server: it must admit an event-log
  abstraction and therefore a conformance surface.
  \item[\textbf{C3}] The phase-transition model of protocol convergence
  (\cref{ch:framework}), including the \emph{latent-heat-of-standardisation}
  formalism that distinguishes incremental integration from discontinuous
  expansion.
  \item[\textbf{C4}] A design-science artifact, \lspmax\ (\cref{ch:artifact}), a
  law-state runtime projected through LSP whose \code{ConformanceVector} keeps
  \admitted, \refused, and \unknownstat\ axes strictly distinct and whose
  diagnostics are emitted as object-centric (OCEL~2.0) events.
  \item[\textbf{C5}] A demonstration (\cref{ch:all-language}) of
  conformance-for-all-language across six domains---natural language, legal text,
  mathematics, biological sequence, music, and business process---unified by the
  object-centric bridge.
  \item[\textbf{C6}] A 2030 roadmap and a sober governance-and-risk analysis
  (\cref{ch:vision}), with every forecast flagged as a \forecast.
\end{itemize}

\section{The Water-to-Steam Metaphor, Made Honest}
\label{sec:metaphor-intro}

Because the metaphor carries argumentative weight, it must be stated with
precision rather than rounded for effect. At $100\,^{\circ}\mathrm{C}$ and one
atmosphere, the specific volume of saturated steam ($\approx 1.673\
\mathrm{m^3/kg}$) divided by that of liquid water ($\approx 0.001044\
\mathrm{m^3/kg}$) gives an expansion factor of about $1{,}603\times$
\citep{engtoolbox-steam,phys-illinois-steam}. The figure of $1{,}700\times$ that
this thesis adopts in its title is the rounded engineering convention used in
boiler-safety practice \citep{nationalboard-steam}; it describes the same
phenomenon at a coarser grain, not a different physical condition. We therefore
speak of a ``${\sim}1{,}600$--$1{,}700\times$'' expansion and treat the precise
ratio as conditions-dependent. This candour is not incidental: the artifact under
study forbids overclaiming, and a thesis about that artifact should model the
discipline it describes.

\Cref{fig:latent-heat} shows the property of phase transitions that the metaphor
turns on. As energy is added to liquid water its temperature rises---until the
boiling point, where temperature \emph{stalls} on a plateau while the substance
absorbs its latent heat and reorganises from liquid to vapour. Only after the
plateau is paid does the expansion appear. Standardisation behaves the same way:
years of specification work register as a flat plateau of apparent
non-progress, after which capability expands discontinuously.

\begin{figure}[t]
\centering
\begin{tikzpicture}[scale=1.05]
  \draw[->,thick] (0,0) -- (10.2,0) node[right]{\footnotesize energy added $Q$ (``standardisation effort'')};
  \draw[->,thick] (0,0) -- (0,5.2) node[above]{\footnotesize state / capability};
  \draw[very thick,lawblue] (0.3,0.5) -- (3,2.7);
  \draw[very thick,refused] (3,2.7) -- (7,2.7);
  \draw[very thick,lawblue] (7,2.7) -- (9.4,4.6);
  \node[lawblue] at (1.5,2.0) {\footnotesize liquid (code only)};
  \node[refused,align=center] at (5,3.15) {\footnotesize latent heat: the protocol is agreed\\\footnotesize (LSP $\to$ MCP $\to$ A2A; LF / AAIF)};
  \node[lawblue,align=center,rotate=38] at (8.5,3.9) {\footnotesize phase change};
  \node[align=center] at (8.7,1.4) {\footnotesize ${\sim}1{,}600$--$1{,}700\times$\\\footnotesize expansion\\\footnotesize (``all language'')};
  \draw[dashed,unknown] (3,0) -- (3,2.7);
  \draw[dashed,unknown] (7,0) -- (7,2.7);
\end{tikzpicture}
\caption[The latent-heat plateau of standardisation]{The latent-heat plateau.
Adding energy raises temperature until the boiling point, where it stalls while
the latent heat of vaporisation is absorbed and the substance expands by
${\sim}1{,}600$--$1{,}700\times$ \citep{nationalboard-steam,latentheat-openstax}.
This thesis reads the standardisation of LSP, MCP, and A2A as the latent heat,
and the generalisation to all language as the phase change.}
\label{fig:latent-heat}
\end{figure}

\section{Structure of the Thesis}
\label{sec:structure}

\Cref{ch:background} surveys the four bodies of work the argument rests on:
process mining and the van der Aalst paradigm; LSP; MCP; and A2A with the wider
agent-interoperability landscape. \Cref{ch:framework} develops the conceptual
framework---``language as process'' and the phase-transition model.
\Cref{ch:methodology} states the design-science methodology and the
bounded-status epistemics inherited from the artifact. \Cref{ch:artifact}
presents \lspmax\ as the artifact. \Cref{ch:fusion} composes LSP, MCP, and A2A
into a single fusion architecture. \Cref{ch:all-language} generalises conformance
checking to six language domains. \Cref{ch:vision} sets out the 2030 vision, the
expansion argument, a roadmap, and risks. \Cref{ch:discussion} revisits the
research questions and the threats to validity, and \cref{ch:conclusion}
concludes. Throughout, claims are graded with bounded status and forecasts are
marked as such.
